% from ega4-1-fr.tex
\oldpage[IV-1]{248}
voisinage ouvert affine $U$ de $X$ ne contenant pas $y$~; comme $X$ est
quasi-séparé, $U$ est rétrocompact dans $X$ (1.2.7), donc $U$ et $X-U$ sont
localement constructibles dans $X$, et par suite ouverts dans
$X^{\mathrm{cons}}$ en vertu de (1.9.11), d'où notre assertion. Comme sur tout
ouvert affine $U$ de $X$, la topologie induite par $\mathscr T^{\mathrm{cons}}$
est celle de $U^{\mathrm{cons}}$ en vertu de (1.9.14, (vi)), il résulte de ce
qui précède que $X^{\mathrm{cons}}$ est localement compact, puisque $U$ est
ouvert dans $X^{\mathrm{cons}}$ et que $U^{\mathrm{cons}}$ est compact. Reste à
prouver que si $X^{\mathrm{cons}}$ est séparé, alors $X$ est quasi-séparé~;
considérons en effet un ouvert affine $U$ de $X$~; la topologie induite sur
$U$ par $\mathscr T^{\mathrm{cons}}$ est celle de $U^{\mathrm{cons}}$
(1.9.14, (vi)), donc $U$ est compact pour cette topologie induite, puisque
$U^{\mathrm{cons}}$ est compact par la première partie du raisonnement. Si
$V$ est un second ouvert affine dans $X$, $U\cap V$ est donc une partie
compacte de l'espace séparé $X^{\mathrm{cons}}$, étant intersection de deux
parties compactes de cet espace~; comme la topologie induite par
$\mathscr T^{\mathrm{cons}}$ sur $U\cap V$ est celle de
$(U\cap V)^{\mathrm{cons}}$ (1.9.14, (vi)) et que l'application identique
$(U\cap V)^{\mathrm{cons}}\to U\cap V$ est continue, on en conclut que
$U\cap V$ est quasi-compact pour la topologie induite par $\mathscr T$~; $X$
est par suite quasi-séparé en vertu de (1.2.7).

(iii) est un corollaire immédiat de (i) et (ii).

(iv) Pour tout $x\in X$, un voisinage ouvert affine $U$ de $x$ pour
$\mathscr T$ est aussi un voisinage de $x$ pour $\mathscr T^{\mathrm{cons}}$
et il est compact en vertu de (iii) et de (1.9.14, (vi)), la topologie induite
sur $U$ par $\mathscr T^{\mathrm{cons}}$ étant identique à celle de
$U^{\mathrm{cons}}$.

(v) Supposons $f$ quasi-compact~; alors on sait déjà (1.9.14, (iv)) que
$f^{\mathrm{cons}}$ est une application fermée. Soit d'autre part $y$ un
point de $Y$, et posons
\[
  Z=f^{-1}(y)=X\times_Y\operatorname{Spec}(k(y));
\]
$Z$ est quasi-compact (1.1.2, (iii)) et comme le morphisme canonique
$p:Z\to X$ est injectif, il résulte de (i) et de ce que l'application
$p^{\mathrm{cons}}$ est continue, que la topologie induite sur $f^{-1}(y)$
par celle de $X^{\mathrm{cons}}$ fait de $f^{-1}(y)$ un espace quasi-compact.
Cela prouve que $f^{\mathrm{cons}}$ est une application propre (Bourbaki,
\emph{Top. gén.}, chap. I\textsuperscript{er}, 3\textsuperscript{e} éd.,
\S~10, n\textsuperscript{o}~2, th.~1). Inversement, supposons que
l'application continue $f^{\mathrm{cons}}$ soit propre, et soit $V$ un ouvert
quasi-compact de $Y$~; si $h:V\to Y$ est l'injection canonique,
$h^{\mathrm{cons}}:V^{\mathrm{cons}}\to Y^{\mathrm{cons}}$ est continue et
injective et $V^{\mathrm{cons}}$ est quasi-compact par (i), donc la topologie
induite sur $V$ par celle de $Y^{\mathrm{cons}}$ fait de $V$ un espace
quasi-compact. L'hypothèse que $f^{\mathrm{cons}}$ est propre entraîne alors
que la topologie induite sur $f^{-1}(V)$ par celle de $X^{\mathrm{cons}}$ fait
de $f^{-1}(V)$ un espace quasi-compact (\emph{loc. cit.}, prop.~6), donc
$f^{-1}(V)$ est aussi un sous-espace quasi-compact de $X$, ce qui montre que
le morphisme $f$ est quasi-compact.
\end{proof}

\begin{env}[1.9.16]
\label{IV.1.9.16-fr}
Nous montrerons plus tard comment on peut, pour tout préschéma $X$, munir
l'espace $X^{\mathrm{cons}}$ d'un faisceau d'anneaux qui en fait un préschéma
dont les anneaux locaux sont tous des corps, identiques aux corps résiduels aux
points de $X$. De tels préschémas s'introduisent par exemple de façon naturelle
dans la traduction, en langage des schémas, des constructions de
Néron [32] dans sa théorie de la réduction des variétés abéliennes.
\end{env}

